%From: Larry Lambe <lambe@matematik.su.se>
%Date: Wed, 15 Mar 1995 11:08:51 +0100
%To: wester@math.unm.edu
%Subject: D. Leites abstract (amsTeX)
%Cc: lambe@insanus.matematik.su.se

\documentstyle{amsppt}
\pageheight{25cm}
\pagewidth{17cm}

\define  \ad{\operatorname{ad}}
\define \adx {(\operatorname{ad} x)}
\define  \rank{\operatorname{rank}}

\define \fg {{\goth{g}}}
\define \fgl {{\goth{gl}}}
\define \fsl {{\goth{sl}}}

\topmatter
\title
Computerizable programs: from supergravity to ballbearings\endtitle

\author P.~Grozman, D.~Leites\endauthor 

\affil Dept\. of Math., Stokholm University\endaffil
\address Roslagsv. 101, 
Kr\"aftriket hus 6,  
S-106 91, Stockholm, Sweden
\endaddress

\thanks Financial support of the Swedish Institute and NFR is greatefully
acknowledged.\endthanks

\keywords cohomology, differential equations, defining relations, gravity, 
Lie algebras, Lie superalgebras, optimal control\endkeywords 

%\subjeclass \endsubjeclass 


\endtopmatter
\document

\subhead 1. Background: supergravity\endsubhead
In 1972, following Berezin's ideas, Leites introduced what is nowadays 
called {\it superschemes}. (Actually, he introduced a more general notion 
pertaining to nonsupercommutaive geometry, cf. Berezin, Shubin, {\it
Schr\"odinger equation},
Kluwer, 1987.) These ideas were not readily shared untill quite recently.
However, after
Wess and Zumino explained (1974) some of physical  advantages of
supersymmetries, the \lq\lq
super" mathematics gained a great deal  of currency. At present a
particular branch of
supermanifold theory, the superstring theory, seem to be an accepted
language of Grand
Unification Theories of fundamental forces (SUSY GUTs). 

In 1970s physicists started to the believe that $N=8$ supergravity is a
good candidate for a
GUT. This theory is free of divergences. This pest of earlier theories
manifests only an
inadequate mathematical description of the physical reality. Still, nobody
could explain what
is the analog of the Riemann tensor for $N$-extended supergravities ---
SUGRA($N$)s  ---
(whereas there are hundreds of papers attempting to do this) since 1974. 

In 70s every physical department had a group working on SUGRA but
eventually only few
entusiasts (or lucky to get grants) continued to persue the topic. 

To write $N=1$ SUGRA equations it took 2 years of trial and error (1974--76). 

The next two steps ($N=2, 3$) were done (with the same  method) by the group of
V.~Ogievetsky in 1984 in JINR, Dubna. A.~Galperin,  E.~Ivanov,
V.~Ogievetsky, E. Sokachev
found out, among other things, that for $N>1$ the original model of the
Minkowski superspace
does not allow one to write SUGRA compatible with  the
Einstein equations on the underlying Minkowski space. (Our results show
that even for $N=1$
the standard model should be reconsidered ...) Even conceding that the
superization of the
Poincar\'e group is correctly found, nobody knows which quotients of the
superPoincar\'e group should we take for THE Minkowski superspace.
 	
Moreover, for $N>3$ the analog of the Riemann tensor is not derived yet in
physical
literature. 

Since it is difficult to seriously discuss for almost 20 years an unwritten
 equation,
though (hypothetically) it is remarkable, the interest to  SUGRA did dry
out everywhere
we must admit. It does not even help that some latest CERN experiments were
interpreted as
the proof of the existence of supersymmetry.

\subhead 2. From supergravity to ballbearings and gyroscopes\endsubhead
Some time ago Leites understood the cause of the difficulty in writing
SUGRA. It is not the 
odd coordinates of the Minkowski superspace (whatever it is for a large $N$). 
The trouble is that for $N>0$ Minkowski superspaces are examples of a 
{\it nonholonomic} supermanifold, i.e., a supermanifold with a nonintegrable 
distribution. 

For nonholonomic manifolds there was no regular method to measure the 
local curvature, such as Riemann tensor for a manifold with a metric. For a
review of
honholonomic problems see [VG] (A.~Vershik,  V.~Gershkovich, Itogi Nauki i
texniki.
Sovr. Probl. matematiki. Fund.  napravleniya, v.18 (Dinamicheskie
sistemy-7) VINITI,1987; in
Russian, to be translated by Springer in Soviet Math. Encyclop. series).
Vershik and 
Gershkovich review numerous problems which go back to Gauss, Carnot, Hertz 
(who coined the term \lq\lq nonholonomic"), Appell, Caratheodory, Schouten, 
et. al.

The simplest example of a nonholonomic system is a ball on a rough plane. More
advanced examples are ballbearings or gyroscopes. The possibility to define an
analog of the Riemann tensor is important in the study of nonholonomic systems:
for instance, it is known that the sign of the curvature is responsible for
stochastisity.

\subhead 3. Optimal control and differential equations\endsubhead
In [VG] and Agrachev's talk at the 1994 Math Congress there were mentioned
problems
of optimal control leading to nonlinear constraints  on generalized
velocities (fields
of cones, spheres, etc.). There are numerous works on (hypoelliptic)
differential equations
(e.g., H\"ormander's) that lead to nonholonomic distributions.

Vershik and Gershkovich gloomily stated that there were no language 
to even formulate the integrability problems for nonholonomic systems 
(with linear constraints to say nothing about nonlinear ones).
	
\subhead 4. How to calculate the local curvature of a nonholonomic
(super)manifold\endsubhead 
Recently we developped a formalism which allows one to 
mathematically rigorously state the integrability problem for any 
nonholonomic system with any constraints. The solution is rather simple,
actually, and
amounts to computing certain Lie (super)algebra cohomology. Being a simple
problem of linear algebra in principle, it is equivalent to an enourmously huge
system of linear  equations. 

In 1992 Grozman wrote a MATHEMATICA based package for calculations of Lie
algebra
(co)homology. He first computed analogues of the Riemann tensor for certain
models of
extended Minkowski superspace. We can now write $N$-SUGRA equations for any
$N=1 - 8$. We
can also corrected calculations of Wess-Zumino constraints from textbooks:
some of the
conditions are redundant.

The huge amount of calculations one has to overcome for $N>1$ does not 
leave any hope to get even a partial answer for large $N$ by trial and
error or by hand
even if we know what to compute and how. (For $N=8$ our program requires
about 255 years to
consider all cases: a day per case.) 


\subhead 5. Formal integrability of partial differential relations\endsubhead
A theorem of S.~Lie states that a symmetry of a differential equation can
be induced from
either 0-jets or from 1-jets of the manifold on which it is considered. For the
first case the criteria of formal integrability are known (Bryan et al.
{\it Exterior
differential systems}, Springer, 1984). 

For the second case our technique gives the answer. We can go even further
and consider not
only equations but certain relations (like those considered by Gromov) as well.

\subhead 6. Noncommutative geometry. (Parastatistics and high temperature
superconductivity)\endsubhead
{\it Super}\-sym\-metry naturally unifies bose and fermi particles. {\it
Metasymmetry}
introduced by Leites and Serganova (Phys. Lett. B, v. 252, n1, 91--96)
naturally describes
particles of anyon type. There is a manifest symmetry wider than
supersymmetry of the 
conventional Lagrangeans if the conventional Minkowski superspaces are
considered as objects
in the category of metaspaces. This seem to be very tempting in connection
with recent
studies of anyons and high temperature superconductivity.
 	
This is also a not too noncommutative generalization of algebraic geometry:
unlike any other
noncommutative geometry, this one preserves all paraphernalia of the usual
differential
geometry: vector fields (hence,  diferential equations), integral, etc.

A further study shows that {\it Volichenko algebras} --- the local versions
of metasymmetry
groups, like Lie algebras ---  introduced and partly classified by Leites
and Serganova
embrace the examples of parasymmetry that appeared in physical papers  of
Rubakov-Spiridonov-Vinet, Bechers-Debergh et al. (Volichenko algebras are
\lq\lq deformations" of 
Lie algebras in a totally new sence, they \lq\lq connect" Lie algebras and
Lie superalgebras.
Their enveloping algebras resemble quantum groups.)

Remarkably, the technique from sec. 4 for computing local invariants can be
applied to a
Minkowski metamanifold  yielding a theory even more general than SUSY GUT,
since it would
describe particles more general than bose-fermi particles.

\subhead 7. Computerizable problems of mathematical physics\endsubhead
It turned out that in order to write a program for calculating the
analogues of Riemannian
tensor mentioned in secs. 2--6 (or other cohomology needed to describe the
structure of Verma
modules, list invariant differentail operators, derive integrability
criteria for
differential equations or relations; or just of a pure mathematical
interest and unswering
some of I.~Gelfand's questions at the 1970 Math. Congress) we have to better
understand the structure of such classical notions as simple Lie algebras
of vector fields.

In particular, we have to calculate their defining relations. (In the end
this leads to 
$q$-quantization of infinite dimensional Lie algebras such as the Poisson
algebra, $\fsl
(\infty)$ and $\fsl(\lambda)$ for a complex $\lambda$.) This is of interest
{\it per se} 
but also in connection with attempts to quantize and $q$-quantize infinite
dimensional 
dynamical systems (like Korteveg-de Vries or Kadomtsev-Petviashvily hierarchies,
etc.) 

Starting with an example of the first complete solution of the {\it
super-Liouville equation}
Leznov and Saveliev taught us a powerful method to formally solve 2-dimensional
equations of mathematical physics. This method requires some knowledge of
representation theory (description of the principal embeddings of $\fsl(2)$,
calculation of Shapovalov's determinant, etc.). Grozman's package allows
one to solve such
problems. We got a presentation of Lie algebras of polynomial vector fields and 
its simple subalgebras. 

\subhead 7.1. Classification of invariant differential operators\endsubhead
This problem is important in applications (see Schouten's 
{\it Tensor calculus for physicists}). All equations of mathematical
physics should
are formulated in invariant terms and certain examples of such 
operators are well known: the exterior differential is the only {\it unary}
invariant operator; Lie derivative, Poisson bracket, Schouten bracket are
examples of {\it binary} operators (classified by Grozman).  To completely
list such operators is {\it Veblen's problem} at 1928 Math. Congress;
Schouten and Niuenhuis returned  to it at later Math Congresses. This
problem can be reduced
to certain Lie  (super)algebra cohomology. This way we discovered plenty of
new operators.

\subhead Conclusion\endsubhead
The above looks as a short summary of major branches of mathematics.
Indeed, Lie algebra
cohomology are ubiquitous. Our package is a set of tools that helps to
calculate many of
these (co)homologies and related problems of linear algebra (Shapovalov
determinant or
defining relations are examples) but to make it really user-friendly (at
least for experts in
MATEMATICA or, for a version, in LISP) is still "a long way to go".


\enddocument
